% ============================================================
%  article_conference.tex
%  Trust-Region vs Line Search — Squelette article conférence
%  Cible : ROADEF 2027 (4-6 p.) ou IFIP TC7 (10-15 p.)
%  Adapter \documentclass selon la venue retenue.
%  Langue : anglais — travail en français.
%
%  Structure (stratégie v2, Section 2.2) :
%    §1  Introduction              (~1 p.)
%    §2  Methods & Protocol        (~1.5 p.)
%    §3  Conditioning Experiment   (~1.5 p.)  ← BLOQUANT G3
%    §4  Per-Category Robustness   (~1.5 p.)
%    §5  Cost & Scalability        (~1 p.)
%    §6  Decision Rule             (~0.75 p.)
%    §7  Theoretical Elements      (~0.75 p.)
%    §8  Conclusion                (~0.5 p.)
%
%  Commandes TODO :
%    \TODO{texte}   — bloc à rédiger
%    \FIG{label}    — placeholder figure
%    \TAB{label}    — placeholder tableau
% ============================================================

% ── Classe de document ──────────────────────────────────────
% ROADEF : remplacer par le template officiel quand disponible
% IFIP TC7 : \documentclass{llncs}  (Springer LNCS)
\documentclass[11pt,a4paper]{article}

% ── Packages ────────────────────────────────────────────────
\usepackage[utf8]{inputenc}
\usepackage[T1]{fontenc}
\usepackage[english]{babel}
\usepackage{amsmath,amssymb,amsthm}
\usepackage{mathtools}           % \coloneqq, etc.
\usepackage{booktabs}            % \toprule, \midrule, \bottomrule
\usepackage{graphicx}
\usepackage{xcolor}
\usepackage{hyperref}
\usepackage[margin=2.5cm]{geometry}
\usepackage{microtype}
\usepackage{natbib}              % \citep, \citet
\usepackage{xspace}
\usepackage{soul}                % \hl{...} pour les TODOs visibles

% ── Commandes utilitaires ────────────────────────────────────
\newcommand{\TODO}[1]{\textcolor{red}{\textbf{[TODO:} \textit{#1}\textbf{]}}}
\newcommand{\FIG}[1]{%
  \begin{center}
    \fbox{\rule{0pt}{3.5cm}\rule{0.85\linewidth}{0pt}}\\[2pt]
    {\small\textit{[Figure placeholder: #1]}}
  \end{center}}
\newcommand{\TAB}[1]{%
  \begin{center}
    \fbox{\parbox{0.85\linewidth}{\centering\vspace{1.2em}
      \small\textit{[Table placeholder: #1]}\vspace{1.2em}}}
  \end{center}}

% ── Notations ───────────────────────────────────────────────
\newcommand{\R}{\mathbb{R}}
\newcommand{\norm}[1]{\lVert#1\rVert}
\newcommand{\Ntot}{N_{\mathrm{tot}}}
\newcommand{\Kmax}{K_{\max}}
\newcommand{\mk}{m_k}
\newcommand{\gk}{g_k}
\newcommand{\Bk}{B_k}
\newcommand{\xk}{x_k}
\newcommand{\dk}{\Delta_k}

% Abréviations méthodes
\newcommand{\SteihaugCG}{Steihaug-CG\xspace}
\newcommand{\NewtonBT}{Newton-BT\xspace}
\newcommand{\BFGSW}{BFGS-Wolfe\xspace}
\newcommand{\BFGSA}{BFGS-Armijo\xspace}

% ============================================================
\begin{document}
% ============================================================

% ── Titre ───────────────────────────────────────────────────
\title{%
  Trust-Region vs Line Search for Unconstrained Optimization:\\
  A Fair-Cost Benchmark on Ill-Conditioned, Non-Convex,\\
  and Indefinite-Hessian Problems%
  \thanks{Code and data: \url{https://github.com/TODO/trust-region-vs-linesearch}
          (tag \texttt{v0.1}).}}

\author{%
  \TODO{Prénom Nom}\\
  Institut National de Statistique et d'Économie Appliquée (INSEA),
  Rabat, Morocco\\
  \texttt{\TODO{email@insea.ac.ma}}}

\date{\TODO{date soumission}}
\maketitle

% ── Abstract ────────────────────────────────────────────────
\begin{abstract}
% Source : stratégie v2, Section 1.3 (abstract proposé).
% À valider/ajuster après G3 (figure pivot κ).
We present a controlled, fair-cost comparison of three trust-region
subproblem solvers (Cauchy point, dogleg, Steihaug--Toint truncated CG)
against three line-search methods (BFGS-Armijo, BFGS-Wolfe, Newton with
backtracking) on 180 unconstrained instances spanning ill-conditioning
($\kappa \in \{10^2, 10^3, 10^4, 10^5\}$), non-convexity, and indefinite
Hessians, with the Hessian-evaluation count separated from the $f/\nabla
f$ budget so that all methods are charged on a common oracle cost
$\Ntot = N_f + N_g + N_h$.
We show that the trust-region family does not dominate line search
globally, and that the operative variable is problem structure rather than
algorithm family.
The \SteihaugCG step emerges as the unique Pareto-efficient choice across
the robustness--scalability plane: it ties the best success rate on
ill-conditioned problems (82.7\%), achieves the best absolute robustness
on non-convex instances (89.1\%), and is the only method preserving a
viable iteration/CPU trade-off from $n=10$ to $n=1000$
($\mathcal{O}(n^2)$ per iteration).
On indefinite Hessians the determining factor is the safeguard mechanism,
not the algorithm family.
We synthesise these findings into a statistically validated decision rule.

\medskip
\noindent\textbf{Keywords:} trust-region methods, line search, unconstrained
optimization, performance profiles, scalability, Steihaug--Toint CG.
\end{abstract}

% ============================================================
\section{Introduction}
\label{sec:intro}
% Matériau : Section 1 du rapport, à réécrire en anglais, resserré.
% Objectif : 1 page, question praticien + contributions + headline chiffré.
% ============================================================

\TODO{%
  Motivation : unconstrained optimization in ML, simulation, calibration.
  Two philosophies : trust-region (simultaneous direction + step via
  $\min_{\norm{p}\le\dk} \mk(p)$) vs line search (decouple direction and step).
  Theoretical guarantees comparable --- practical behavior diverges.
  Central question: for which problem structures does each family excel?
}

\TODO{%
  Headline figure : Steihaug-CG 82.7\% on Cat.~A, 89.1\% on Cat.~B;
  Newton-BT 76\% on Cat.~C; global averages TR 62.6\% vs LS 65.7\%
  (no global dominance).
}

\paragraph{Contributions.}
\begin{enumerate}
  \item A complete, from-scratch implementation of six second-order methods
        with a common oracle-cost framework separating $N_f$, $N_g$, $N_h$.
  \item A structured benchmark of 180 instances in three difficulty
        categories (ill-conditioning, non-convexity, indefinite Hessian)
        across dimensions $n \in \{10, 50, 100, 500, 1000\}$.
  \item \TODO{Quantified $\kappa$-resolved experiment (§\ref{sec:cond})
        --- fill after G3.}
  \item A statistically validated decision rule synthesising all findings
        (§\ref{sec:rule}).
\end{enumerate}

% ============================================================
\section{Methods and Experimental Protocol}
\label{sec:methods}
% Matériau : §§3--4 du rapport.
% Objectif : 1.5 pages, tableau compact 6 méthodes, protocole budget unifié.
% ============================================================

\subsection{Trust-Region Methods}

At iteration $k$, solve the subproblem
$\min_{\norm{p} \le \dk} \mk(p)$
where $\mk(p) = f_k + \gk^\top p + \tfrac{1}{2} p^\top \Bk p$.
The trust-region radius $\dk$ is updated via the ratio
$\rho_k = \mathrm{ared}_k / \mathrm{pred}_k$
with thresholds $\eta_1 = 0.1$, $\eta_2 = 0.75$.
Three subproblem solvers are compared:

\begin{description}
  \item[Cauchy Point (CP).] Minimizes $\mk$ along $-\gk$;
        $\tau^* = \min(1, \norm{\gk}^3 / (\gk^\top\Bk\gk \cdot \dk))$.
  \item[Dogleg.] Interpolates the Cauchy direction $p^U$ and the Newton
        step $p^N = -\Bk^{-1}\gk$ (Cholesky; fallback to CP if $\Bk \not\succ 0$).
  \item[\SteihaugCG.] Truncated CG solving $\Bk p = -\gk$; stops on
        negative curvature ($d^\top\Bk d \le 0$), boundary, or
        Eisenstat--Walker tolerance \citep{Steihaug83}.
\end{description}

\subsection{Line-Search Methods}

\TODO{BFGS-Armijo, BFGS-Wolfe (zoom + cubic interpolation), Newton-BT
(fallback $-g_k$ when $H_k$ singular; §\ref{sec:cat_c} discusses the
implication for Cat.~C).}

\subsection{Common Protocol}

\TAB{Table: 6 methods --- compact parameter table.
Rows: method name, line/TR, Kmax=1000 (uniform, G8), tol=1e-6,
oracle cost Ntot=Nf+Ng+Nh.}

\paragraph{Oracle-cost budget (G8).}
A uniform iteration budget $\Kmax = 1000$ is applied to all six solvers.
The total oracle cost $\Ntot = N_f + N_g + N_h$ is tracked separately for
each method; methods that never call the Hessian oracle (BFGS variants)
naturally incur $N_h = 0$, making comparisons equitable.
Performance is reported via Dolan--Moré performance profiles
\citep{DolanMore2002} and data profiles \citep{MoreWild2009}.

\subsection{Benchmark Structure}

\TAB{Table: benchmark --- 3 categories × 5 dimensions × nb problems
= 180 instances. Cat.~A: quadratics with $\kappa \in \{10^2,...,10^5\}$,
Rosenbrock, Wood, Dixon-Price; Cat.~B: non-convex (Rastrigin, Ackley,
CUTEst BROYDN3D, CHAINWOO, ...); Cat.~C: indefinite Hessian at $x^*$.}

% ============================================================
\section{The Conditioning Experiment: Central Result}
\label{sec:cond}
% *** BLOQUANT G3 ***
% Matériau : À produire (gap G3).
% Objectif : 1.5 pages, Figure 1 pivot + narration centrale.
% Milestone : Steihaug ≈ constant sur κ ∈ [1e2,1e5], BFGS croissant.
%             Si découplage absent → pause + rediscussion du récit.
% ============================================================

\TODO{%
  SECTION BLOQUANTE --- attend gap G3 (figure itérations médianes vs $\kappa$).\\
  Structure prévue :\\
  (i)  Experimental design: Cat.~A quadratics at fixed $n$,
       $\kappa \in \{10^2, 10^3, 10^4, 10^5\}$ varied.\\
  (ii) Figure~\ref{fig:kappa}: median iterations vs $\kappa$ for
       Steihaug-CG, Dogleg, BFGS-Armijo, BFGS-Wolfe.\\
  (iii) Companion panel: vs $n$ at fixed $\kappa$.\\
  (iv) Wilcoxon signed-rank test: Steihaug vs BFGS-Armijo intra-Cat.~A
       (optional at conference stage, full stats for journal).\\
  Expected finding: Steihaug median iterations $\approx 21$ regardless
  of $\kappa$; BFGS-Armijo rises from $\sim$20 to $\sim$200 as
  $\kappa: 10^2 \to 10^5$.
}

\begin{figure}[htbp]
  \centering
  \FIG{fig:kappa --- PIVOT FIGURE (G3). Left panel: median iterations
       vs $\kappa$ at fixed $n$. Right panel: vs $n$ at fixed $\kappa$.
       4 methods: Steihaug-CG (flat), Dogleg (flat), BFGS-Armijo
       (rising), BFGS-Wolfe (rising).}
  \caption{Median iterations as a function of conditioning $\kappa$
    (left) and dimension $n$ (right) on Category~A problems.
    \SteihaugCG and Dogleg maintain a near-constant iteration count
    while BFGS variants degrade with $\kappa$.
    \TODO{Fill after G3.}}
  \label{fig:kappa}
\end{figure}

\TODO{Explain the mechanism: TR methods adapt $\dk$ to local curvature
independently of $\kappa$; BFGS accumulates curvature information
progressively and degrades on highly elongated landscapes.
Theoretical backing: CG convergence rate $1 - 1/\kappa$ \citep[Ch.~5]{NW06};
Cauchy bound $\tfrac{1}{2}\norm{\gk}\min(\dk, \norm{\gk}/\norm{\Bk})$
\citep[Ch.~4]{CGT00}.}

% ============================================================
\section{Per-Category Robustness}
\label{sec:robustness}
% Matériau : §§5.2 + 5.4 du rapport, mis à jour post-G1 (Newton-LM).
% Objectif : 1.5 pages.
% ============================================================

\TAB{Table: success rates (\%) by category and method.
Rows: 6 methods. Columns: Cat.~A, Cat.~B, Cat.~C, Global.
(Values from rapport Table 5 --- update after re-run v2.)}

\paragraph{Category~A (ill-conditioned).}
\TODO{Dogleg and Steihaug-CG dominate (82.7\%) vs BFGS (62--68\%).
Cauchy Point collapses (10.7\%): steepest-descent direction is
ineffective in elongated valleys.
Reference Fig.~\ref{fig:kappa}.}

\paragraph{Category~B (non-convex).}
\TODO{\SteihaugCG achieves 89.1\%, best absolute score, due to negative
curvature detection ($d^\top\Bk d \le 0$ exit with boundary projection).
BFGS-Armijo competitive at 74.5\%.}

\paragraph{Category~C (indefinite Hessian).}
\label{sec:cat_c}
\TODO{The operative variable is the \emph{safeguard mechanism}, not the
algorithm family (Q1-acté finding).
Newton-BT survives at 76\% via its $-g_k$ fallback (ad-hoc safeguard).
BFGS-Wolfe collapses at 26\% (26 numerical exceptions, no safeguard).
\SteihaugCG reaches 58\% without code modification via its
$d^\top\Bk d \le 0$ criterion.
\textbf{G1 note:} after Newton-LM regularisation, Newton-BT Cat.~C
results should be updated here.}

\paragraph{Performance and data profiles.}

\begin{figure}[htbp]
  \centering
  \FIG{fig:profiles --- Dolan-Moré performance profiles (3 panels,
       one per category) + data profiles (Moré \& Wild 2009).
       Matériau : rapport Figs 6--8 (Dolan-Moré) + data\_profiles\_phase0.pdf (G8).}
  \caption{Performance profiles (Dolan--Moré) and data profiles
    (Moré--Wild) per category.
    \TODO{Insert after re-run v2 and data profiles script (G8).}}
  \label{fig:profiles}
\end{figure}

% ============================================================
\section{Cost and Scalability}
\label{sec:scalability}
% Matériau : §§5.3 + 5.5 du rapport.
% Objectif : 1 page.
% ============================================================

\paragraph{Oracle cost vs dimension.}
\TODO{Fig: $\Ntot$/CPU vs $n$ for 3 categories.
Newton-BT: $\mathcal{O}(n^3)$ per iteration, prohibitive at $n=1000$.
\SteihaugCG: $\mathcal{O}(n^2)$ per iteration, stable cost.
Message: Steihaug is the only method maintaining a reasonable
iteration/CPU trade-off from $n=10$ to $n=1000$.}

\begin{figure}[htbp]
  \centering
  \FIG{fig:scalability --- $\Ntot$ median (left) and CPU (right) vs $n$.
       Matériau : rapport Fig.~5 (updated post-v2).}
  \caption{Median oracle cost $\Ntot$ and CPU time as a function of $n$.
    \TODO{Fill after re-run v2.}}
  \label{fig:scalability}
\end{figure}

\paragraph{Rosenbrock case study.}

\TAB{Table: Rosenbrock $n=10$, $x_0=(-1.2,1,\ldots)^\top$.
3 methods: Steihaug-CG, BFGS-Wolfe, Newton-BT.
Columns: iterations, $\norm{\nabla f}_{\mathrm{final}}$, $\#f$, $\#\nabla f$.
(Rapport Table 6.)}

\TODO{Steihaug-CG: 39 iterations, 77 $f$-evals --- best cost/accuracy.
Newton-BT: 291 iterations, 2602 $f$-evals (frequent backtracking rejections).}

% ============================================================
\section{Decision Rule}
\label{sec:rule}
% Matériau : §7.2 du rapport, formalisé en arbre if/then.
% Objectif : 0.75 page --- À produire (gap G2).
% ============================================================

\TODO{%
  Figure: decision tree with 4 nodes:\\
  (1) Is $\nabla^2 f$ available?\\
  (2) Is $H$ potentially indefinite?\\
  (3) Is $\kappa$ large (ill-conditioned)?\\
  (4) Is $n$ large ($\ge 500$)?\\
  Leaves: Steihaug-CG, Dogleg, Newton-BT (small $n$), BFGS-Armijo.\\
  At conference stage: qualitative (no statistical test).\\
  At journal stage: backed by Wilcoxon signed-rank + Holm correction (G7).
}

\begin{figure}[htbp]
  \centering
  \FIG{fig:decision-tree --- if/then tree (4 nodes, 4-5 leaves).
       Gap G2.}
  \caption{Recommended solver selection rule based on problem structure.
    \TODO{Produce after G2.}}
  \label{fig:decision}
\end{figure}

% ============================================================
\section{Theoretical Elements}
\label{sec:theory}
% Matériau : §6 du rapport.
% Objectif : 0.75 page --- esquisse des preuves.
% ============================================================

\subsection{Cauchy Point Bound}

\begin{proposition}[Sufficient reduction, \citealt{CGT00}]
\label{prop:cauchy}
The Cauchy step $p^C_k$ satisfies
\begin{equation}
  \mk(0) - \mk(p^C_k)
  \;\ge\;
  \tfrac{1}{2}\norm{\gk}
  \min\!\left(\dk,\,\frac{\norm{\gk}}{\norm{\Bk}}\right).
  \label{eq:cauchy_bound}
\end{equation}
\end{proposition}

\begin{proof}[Proof sketch]
\TODO{Two cases: $\gk^\top\Bk\gk \le 0$ and $> 0$.
Case 1 ($\gk^\top\Bk\gk \le 0$): $\alpha = \dk/\norm{\gk}$,
reduction $= \dk\norm{\gk} \ge \tfrac{1}{2}\norm{\gk}\dk$. $\checkmark$\\
Case 2 ($\gk^\top\Bk\gk > 0$): $\alpha^* = \norm{\gk}^2/(\gk^\top\Bk\gk)$;
distinguish $\alpha^* \le \dk/\norm{\gk}$ vs $> \dk/\norm{\gk}$;
assemble the two bounds. $\checkmark$}
\end{proof}

\subsection{Global Convergence}

\begin{theorem}[Global convergence, \citealt{CGT00}]
If $f \in C^2(\R^n)$ is bounded below and $\norm{\Bk}$ is bounded, then any
algorithm ensuring reduction \eqref{eq:cauchy_bound} satisfies
$\liminf_{k\to\infty} \norm{\nabla f(\xk)} = 0$.
\end{theorem}

\begin{proof}[Proof sketch]
\TODO{By contradiction: if $\norm{\gk} \ge \varepsilon > 0$, then
$\mathrm{pred}_k \ge c\varepsilon^2 > 0$.
For $\dk$ small, $\rho_k \to 1$ (Taylor; requires $\nabla f$
$L$-Lipschitz), so $\mathrm{ared}_k \ge \eta_1\,\mathrm{pred}_k > 0$,
$\sum_k \mathrm{ared}_k$ diverges, contradicting boundedness below.
(Full argument with Lipschitz step: journal version, §7.)}
\end{proof}

\subsection{Asymptotic Rates}

\TODO{CP: linear at rate $1-1/\kappa(\Bk)$ (explains 10.7\% Cat.~A).
Dogleg/Steihaug with $\Bk = \nabla^2 f(\xk)$: superlinear.
Newton-BT: quadratic local, $\norm{x_{k+1}-x^*} = \mathcal{O}(\norm{x_k-x^*}^2)$
(explains 9 median iterations).}

% ============================================================
\section{Conclusion}
\label{sec:conclusion}
% Matériau : §8 du rapport.
% Objectif : 0.5 page.
% ============================================================

\TODO{%
  Summary per category (3 bullet points from rapport §8).\\
  Main finding: Steihaug-CG is the unique Pareto-efficient solver across
  the robustness-scalability plane (Q1-acté).\\
  Cat.~C: safeguard mechanism, not family (Q1-acté).\\
  Limitations: hand-reimplemented CUTEst problems (pycutest for journal);
  no L-BFGS comparison (journal, G4); dense Hessians only (sparse: journal G5).\\
  Perspectives: hybrid TR/quasi-Newton; L-BFGS extension; pycutest.
}

% ── Acknowledgements ────────────────────────────────────────
\section*{Acknowledgements}
\TODO{INSEA context --- à adapter.}

% ── Bibliography ────────────────────────────────────────────
\bibliographystyle{plainnat}
\bibliography{references_conf}

\end{document}

% ============================================================
% END OF article_conference.tex
% ============================================================
